% ЗГЕНЕРОВАНО scripts/build_format_v2_chapter.py — не редагувати вручну.
% Джерело: books/ashtavakra-gita/poetychnyi/rozdil-20.md

\poemverse{20.1}{Де першоелементи, де тіло, \\ де чуття, де розум? \\ Де порожнеча, де без-надійність — \\ у моїй власній формі, \\ що не знає плями?}
\poemverse{20.2}{Де писання, де пізнання Я, \\ де розум без обʼєктів? \\ Де вдоволеність, де свобода від спраги — \\ у мене, в кого пари вже відпали назавжди?}
\poemverse{20.3}{Де знання й де незнання, \\ де я, де це, де моє? \\ Де пута, де визволення — \\ і де взагалі якась оформленість \\ у того, хто сам є своєю формою?}
\poemverse{20.4}{Де вже запущені в дію діяння, \\ де звільнення ще за життя? \\ Де те безтілесне кайвальям — \\ у того, хто завжди без ознак?}
\poemverse{20.5}{Де діяч, де той, хто зазнає? \\ Де бездіяльне сяйво? \\ Де безпосереднє знання, де його плід — \\ у мене, в кого немає власної природи?}
\poemverse{20.6}{Де світ, де той, хто прагне звільнення, \\ де йог чи знавець? \\ Де звʼязаний, де звільнений — \\ я у своїй власній формі, \\ яка не знає двоїстості.}
\poemverse{20.7}{Де творення, де згортання, \\ де мета, де засіб? \\ Де той, хто практикує, де здобуття — \\ я у своїй власній формі, \\ яка не знає двоїстості.}
\poemverse{20.8}{Де той, хто пізнає, де засіб пізнання, \\ де предмет, де саме пізнання? \\ Де щось, де ніщо — \\ у мене, хто завжди без плями?}
\poemverse{20.9}{Де розсіяність, де зосередженість, \\ де несвідомість, де затьмареність? \\ Де радість, де відчай — \\ у мене, хто завжди бездіяльний?}
\poemverse{20.10}{Де ця буденна практика, \\ де та вища істина? \\ Де щастя, де страждання — \\ у мене, хто вже завжди нічого не розглядає?}
\poemverse{20.11}{Де мая, де сансара, \\ де прихильність і де відвернення? \\ Де жива істота, де те Брахмо — \\ у мене, хто завжди без плями?}
\poemverse{20.12}{Де діяння і де припинення, \\ де звільнення, де пута — \\ у мене, незмінного й неділимого, \\ що завжди в собі?}
\poemverse{20.13}{Де настановлення, де писання, \\ де учень, де вчитель? \\ Де й мета людського життя — \\ у мене, хто без жодної опори, \\ у мене, блаженного?}
\poemverse{20.14}{Де є і де немає, \\ де одне і де двоє? \\ Та й навіщо багато слів: \\ у мене не постає ніщо.}
